% NOETHER PAPER 6 — KOREAN TRANSLATION-PRODUCER DRAFT — UNCHECKED
% Unit: T05-U45; current-authority whole lines 4976--4989
% Source slice: 684 LF UTF-8 bytes; SHA-256 0E8CE26842AA803C24EB0227F14F6144CA0F040DE774F187823FB2FF60EDCA89
% Pointer: NOETH-DE-AUTH-v007-20260804; SHA-256 A6A8FC8E5AC24ACAF49DFD55B4B58FA3DA882EF8C3FDD4D136220C8751045156
% German authority: NOETH-DE-ED-0001; 2,153,565 raw bytes; SHA-256 D1F06B311F6CBD991DD247D745DD9A72DDE326A20396DF43CFE0C8EDB1593CDB
% Producer translation only: no source/Korean/formula review, compilation, or rendering.
% Sense window: birationale Transformation is the algebraic-geometric field-generator change, not a generic reversible map.
% Terminology debt: birationale Transformation→쌍유리 변환/쌍유리변환 and Rationalbasis→유리기저/유리 기저 remain checker-held.
% Standardization hold: ko-KR 쌍유리 변환 is provisional; ko-KP equivalent and any Hanja supplement remain unresolved.

$\Kfield_{n\rho}$에서 택한 $\rho$개의 대수적으로 독립인 함수로 이루어진 임의의 계가 $h_1(x)\cdots h_\rho(x)$로 주어졌다고 하자. 이 계는 위의 방법으로 다음 유리기저까지 확장할 수 있다.
\[
h_1(x)\cdots h_\rho(x);\ h_{\rho+1}(x)\cdots h_\sigma(x)
\]
정의 성질에 따라 다음 관계식들이 존재한다.
\[
\begin{array}{ll}
h_i(x)=\chi_i\bigl(g_1(x)\cdots g_k(x)\bigr)&(i=1,2,\ldots,\sigma),\\
g_j(x)=\psi_j\bigl(h_1(x)\cdots h_\sigma(x)\bigr)&(j=1,2,\ldots,k),
\end{array}
\]
여기서 $\chi_i$, $\psi_j$는 각각 그 인수들의 유리함수를 뜻한다. 즉,

임의의 한 유리기저로부터 다른 모든 유리기저가 쌍유리 변환으로 얻어진다. 또한,
